\chapitre{Comparaison des décompositions OpenMP à grande largeur}
Le nombre de synchronisations et la largeur des étapes distinguent les deux décompositions. Pour une matrice de \(m\times n\) cases, les anti-diagonales ont \(m+n-1\) niveaux intérieurs de largeur variable. NvPD traite \(m\) lignes de largeur \(n\) et lit un prédécesseur indirect dans l'index des occurrences.

La comparaison OpenMP reprend les quinze entrées d'un million de bases avec l'exécutable de la campagne initiale, antérieure à la vectorisation. Chaque version emploie huit cœurs. À chaque vague, trois processus par nœud traitent les trois taux sur des cœurs distincts. Les nœuds échangent les algorithmes entre deux vagues. Les cinq valeurs d'un taux correspondent à cinq paires différentes.
\chart{13_openmp_un_million.png}{Les durées OpenMP sur quinze entrées et le gain SIMD sur trois de ces entrées proviennent de deux campagnes distinctes à huit cœurs par processus.}{fig:million}{bench/complete_1m.sbatch}
Sur le panneau a) de la figure~\ref{fig:million}, les médianes des anti-diagonales vont de 190,04 à 190,29\,s selon le taux. Les médianes de NvPD vont de 309,96 à 310,91\,s. Les trente scores concordent avec les deux bibliothèques.
\synthese{À 5\,\% d'éditions injectées, la version initiale de NvPD prend 1,63 fois le temps des anti-diagonales.}{Le panneau b) isole l'effet de la vectorisation de sa boucle de colonnes.}
